\documentclass[12pt,a4paper]{article}

% -------------------------------------------------
% Paquetes y configuración
% -------------------------------------------------
\usepackage[spanish,es-nodecimaldot]{babel}
\usepackage[utf8]{inputenc}
\usepackage[T1]{fontenc}
\usepackage{lmodern}
\usepackage[a4paper,margin=2.5cm]{geometry}
\usepackage{enumitem}
\usepackage{longtable}
\usepackage{array}
\usepackage{booktabs}
\usepackage{xcolor}
\usepackage{hyperref}
\usepackage{titlesec}
\usepackage{fancyhdr}
\usepackage{microtype}

\definecolor{azulhotel}{HTML}{17365D}
\definecolor{azulclaro}{HTML}{D9EAF7}
\definecolor{grisclaro}{HTML}{F2F2F2}
\hypersetup{
  colorlinks=true,
  linkcolor=azulhotel,
  urlcolor=azulhotel,
  pdftitle={Requisitos e historias de usuario - Sistema de Gestión Hotelera Integral},
  pdfauthor={Documento de análisis de software}
}
\titleformat{\section}{\Large\bfseries\color{azulhotel}}{\thesection.}{0.5em}{}
\titleformat{\subsection}{\large\bfseries\color{azulhotel}}{\thesubsection}{0.5em}{}
\pagestyle{fancy}
\fancyhf{}
\lhead{\small Sistema de Gestión Hotelera Integral}
\rhead{\small Requisitos de software}
\cfoot{\thepage}
\setlength{\headheight}{14.5pt}
\setlist[itemize]{itemsep=2pt,topsep=3pt}
\setlist[enumerate]{itemsep=3pt,topsep=3pt}
\renewcommand{\arraystretch}{1.2}

\newcommand{\RF}[2]{\item[\textbf{#1}] #2}
\newcommand{\RNF}[2]{\item[\textbf{#1}] #2}
\newcommand{\HU}[5]{%
  \subsection*{#1 -- #2}
  \noindent\textbf{Prioridad:} #3 \hfill \textbf{Fase:} #4\\
  \noindent\textbf{Historia:} #5
  \vspace{2mm}
}

\begin{document}

\begin{titlepage}
  \centering
  \vspace*{2cm}
  {\Large\bfseries DOCUMENTO DE REQUISITOS DE SOFTWARE\par}
  \vspace{0.8cm}
  {\huge\bfseries Sistema de Gestión Hotelera Integral\par}
  \vspace{0.3cm}
  {\Large Plataforma multitenant para cuatro hoteles asociados\par}
  \vspace{1.5cm}
  \color{azulhotel}\rule{0.85\textwidth}{1.2pt}
  \vspace{1cm}

  \begin{tabular}{>{\bfseries}p{4cm} p{8cm}}
    Producto & Sistema de gestión hotelera integral (PMS)\\
    Organización & Cadena hotelera asociada, compuesta por cuatro sedes\\
    Usuarios destinatarios & Junta de socios, administradores, recepción, camarería y auditoría\\
    Backend previsto & .NET 9 o superior, ASP.NET Core Web API\\
    Frontend previsto & Angular\\
    Base de datos preferida & SQL Server\\
    Documento & Requisitos funcionales, no funcionales e historias de usuario\\
  \end{tabular}

  \vfill
  {\large Documento elaborado a partir de la especificación del proyecto\par}
\end{titlepage}

\tableofcontents
\newpage

\section{Introducción}

Este documento organiza los requisitos funcionales, los requisitos no funcionales y las historias de usuario del proyecto \textit{Sistema de Gestión Hotelera Integral Multitenant}. La solución busca centralizar la operación de cuatro hoteles asociados, manteniendo el aislamiento lógico de la información de cada sede y ofreciendo a la Junta de Socios una visión consolidada de la cadena.

Actualmente, las reservas, los pagos, los consumos y los informes se administran mediante archivos de Excel y procesos manuales. Esto genera demoras, duplicidad de reservas, falta de trazabilidad, errores de cobro y dificultad para consultar indicadores en tiempo real.

\section{Objetivo y alcance}

\subsection{Objetivo general}
Definir las capacidades y condiciones de calidad que debe cumplir una plataforma hotelera centralizada para gestionar reservas, huéspedes, habitaciones, tarifas, consumos, pagos, limpieza, reportes y auditoría.

\subsection{Alcance funcional}
El sistema incluirá:
\begin{itemize}
  \item Gestión multihotel y aislamiento de datos por sede.
  \item Autenticación, usuarios, roles, permisos y asignación de hoteles.
  \item Parametrización de hoteles, habitaciones, tarifas, temporadas, impuestos, políticas y servicios.
  \item Cotización, disponibilidad, reservas, cancelaciones, check-in, consumos, pagos y check-out.
  \item Gestión básica de limpieza de habitaciones.
  \item Notificaciones por correo y comprobantes digitales.
  \item Dashboard consolidado, reportes exportables y auditoría.
  \item Consentimiento para tratamiento de datos personales y anonimización según permisos.
\end{itemize}

\subsection{Fuera del alcance}
La pasarela de pagos en línea no forma parte de esta versión. Los pagos se registrarán manualmente indicando método y referencia. El portal de autogestión del huésped se contempla como una funcionalidad de fase posterior (HU-24).

\section{Actores y roles}

\begin{longtable}{p{3.2cm} p{10.2cm}}
\toprule
\textbf{Actor} & \textbf{Responsabilidades principales}\\
\midrule
Socio & Consulta de solo lectura de los indicadores consolidados de la cadena y de cada sede.\\
Administrador Global & Configura hoteles, parámetros globales, empleados y roles; consulta información consolidada y gestiona procesos autorizados.\\
Administrador de Sede & Administra la operación de su hotel, incluyendo habitaciones, reservas, personal y parámetros permitidos.\\
Recepcionista & Consulta disponibilidad, crea y actualiza reservas, realiza check-in/check-out, registra consumos y pagos.\\
Camarería / Supervisor & Consulta habitaciones asignadas y actualiza los estados de limpieza.\\
Auditor & Consulta información de auditoría y exporta reportes de acuerdo con sus permisos.\\
Huésped & Proporciona sus datos y consentimiento; recibe confirmaciones, vouchers y comprobantes. En esta versión no inicia sesión para crear reservas.\\
\bottomrule
\end{longtable}

Los usuarios de sede operan únicamente dentro del hotel activo asignado. Los permisos deben aplicarse en el servidor y no depender solamente de la interfaz.

\section{Requisitos funcionales}

Los siguientes requisitos se derivan de las capacidades y reglas descritas en el documento fuente. La numeración RF permite referenciarlos durante el diseño, desarrollo y pruebas.

\subsection{Acceso, usuarios y permisos}
\begin{description}[leftmargin=2.1cm,style=nextline]
  \RF{RF-01}{El sistema debe permitir iniciar sesión mediante correo electrónico y contraseña.}
  \RF{RF-02}{El sistema debe administrar sesiones mediante un token de acceso de corta duración y un mecanismo de renovación seguro, con revocación y detección de reutilización de tokens.}
  \RF{RF-03}{El sistema debe bloquear temporalmente la cuenta después de cinco intentos fallidos consecutivos y registrar el evento en auditoría.}
  \RF{RF-04}{El Administrador Global debe poder crear, modificar, desactivar y asignar empleados a roles, departamentos, turnos y hoteles.}
  \RF{RF-05}{Si un empleado tiene asignados varios hoteles, el sistema debe permitirle seleccionar como contexto activo únicamente una sede autorizada.}
  \RF{RF-06}{El sistema debe aplicar permisos granulares por rol y acción. Una operación no autorizada dentro de la sede propia debe responder con error 403.}
  \RF{RF-07}{El sistema debe impedir que un usuario consulte o modifique recursos operativos de otra sede. Los recursos pertenecientes a otra sede deben responder como no encontrados (404).}
\end{description}

\subsection{Parametrización de hoteles y servicios}
\begin{description}[leftmargin=2.1cm,style=nextline]
  \RF{RF-08}{El Administrador Global debe poder crear y configurar hoteles, incluyendo nombre, ciudad, dirección, zona horaria, moneda e identidad visual.}
  \RF{RF-09}{El sistema debe impedir desactivar un hotel que tenga reservas futuras activas e informar la cantidad de reservas que lo impiden.}
  \RF{RF-10}{El sistema debe permitir definir tipos de habitación, capacidad, comodidades y habitaciones físicas por sede.}
  \RF{RF-11}{El sistema debe impedir números de habitación duplicados dentro de un mismo hotel, pero permitir que dos hoteles distintos utilicen el mismo número.}
  \RF{RF-12}{El sistema debe impedir desactivar habitaciones que tengan reservas futuras activas.}
  \RF{RF-13}{El sistema debe permitir configurar tarifas base, temporadas, recargos de fin de semana y promociones por tipo de habitación.}
  \RF{RF-14}{El sistema debe rechazar temporadas con fechas traslapadas dentro de la misma configuración hotelera.}
  \RF{RF-15}{El sistema debe aplicar promociones vigentes y mostrar el descuento en el desglose de la tarifa.}
  \RF{RF-16}{El sistema debe permitir configurar impuestos y políticas de cancelación por sede. El cálculo de cancelación debe usar la política vigente al momento de realizar la reserva.}
  \RF{RF-17}{El Administrador de Sede debe poder administrar un catálogo de servicios adicionales con precio por hotel, por ejemplo desayuno, spa, parqueadero, minibar y lavandería. Los servicios inactivos no deben estar disponibles para registrar consumos.}
\end{description}

\subsection{Disponibilidad y reservas}
\begin{description}[leftmargin=2.1cm,style=nextline]
  \RF{RF-18}{El recepcionista debe poder consultar disponibilidad por rango de fechas, tipo de habitación y cantidad de huéspedes, con precio total y desglose por noche.}
  \RF{RF-19}{La consulta de disponibilidad no debe mostrar habitaciones ocupadas por reservas activas que se traslapen con el rango solicitado.}
  \RF{RF-20}{El sistema debe permitir consultar disponibilidad entre sedes y devolver únicamente sede, tipo de habitación, cantidad disponible y tarifa referencial. Esta consulta debe quedar auditada.}
  \RF{RF-21}{La consulta entre sedes no debe exponer reservas, datos de huéspedes, números de habitación ni consumos de otras sedes.}
  \RF{RF-22}{El recepcionista debe poder crear una reserva en un flujo integrado con huésped titular, fechas, habitación, tarifa, descuentos e impuestos. La reserva inicial debe quedar en estado \textit{Pendiente}.}
  \RF{RF-23}{El sistema debe admitir una clave de idempotencia al crear reservas, de forma que repetir la misma petición no genere una reserva duplicada.}
  \RF{RF-24}{El sistema debe impedir el overbooking mediante una garantía de integridad en la base de datos para evitar reservas activas de una misma habitación con fechas traslapadas.}
  \RF{RF-25}{El sistema debe detectar ediciones concurrentes de una misma reserva o habitación mediante control de concurrencia optimista y comunicar el conflicto.}
  \RF{RF-26}{El sistema debe validar las transiciones de estado de las reservas: Pendiente, Confirmada, Check-In, Check-Out, Liquidada, Anulada y No-Show, de acuerdo con las reglas de negocio.}
  \RF{RF-27}{El sistema debe registrar en el historial de cada transición el estado anterior, el nuevo estado, el usuario, la fecha y el motivo cuando corresponda.}
  \RF{RF-28}{El sistema debe permitir cancelar reservas o registrar No-Show de acuerdo con la política de la sede, calcular la penalidad aplicable y liberar la disponibilidad cuando corresponda.}
  \RF{RF-29}{El sistema debe ofrecer un calendario de ocupación por habitación, con un rango de consulta de hasta 14 días y estados identificables por color y texto.}
\end{description}

\subsection{Recepción, estancia, consumos y pagos}
\begin{description}[leftmargin=2.1cm,style=nextline]
  \RF{RF-30}{El recepcionista debe poder registrar al huésped titular y a sus acompañantes, incluyendo los datos de identificación requeridos para la estancia.}
  \RF{RF-31}{El sistema debe impedir completar el check-in si la habitación no está en estado Limpia/Disponible o si no se ha registrado el consentimiento de tratamiento de datos requerido.}
  \RF{RF-32}{Al completar el check-in, la reserva debe pasar a estado Check-In y la habitación a Ocupada.}
  \RF{RF-33}{El sistema debe permitir que camarería consulte y actualice el estado de limpieza de las habitaciones desde una interfaz adaptable a tabletas.}
  \RF{RF-34}{Al completar el check-out, la habitación debe pasar automáticamente al estado Sucia.}
  \RF{RF-35}{El recepcionista debe poder registrar consumos y cargos en el folio del huésped, incluyendo cantidad, valor, usuario y hora.}
  \RF{RF-36}{Los cargos ya registrados no deben editarse directamente; su corrección debe realizarse mediante anulación con motivo y contra-asiento auditado.}
  \RF{RF-37}{El folio debe mostrar cargos, impuestos, pagos y saldo pendiente.}
  \RF{RF-38}{El sistema debe permitir registrar pagos y abonos con método y referencia, actualizar el saldo y generar un recibo.}
  \RF{RF-39}{El registro de pagos debe admitir una clave de idempotencia para evitar duplicaciones por reintentos de red.}
  \RF{RF-40}{El sistema debe impedir el check-out cuando el saldo del folio sea distinto de cero y debe informar el saldo pendiente.}
  \RF{RF-41}{Cuando el saldo sea cero, el sistema debe permitir cerrar la estancia, generar un comprobante PDF y pasar la reserva a Liquidada cuando el folio esté cerrado y el comprobante final haya sido emitido.}
  \RF{RF-42}{El sistema debe permitir cerrar la caja por turno y congelar los totales por método de pago del cierre realizado.}
\end{description}

\subsection{Notificaciones, privacidad y datos personales}
\begin{description}[leftmargin=2.1cm,style=nextline]
  \RF{RF-43}{El sistema debe enviar correos transaccionales con el detalle de la reserva, voucher con código QR y comprobantes, utilizando la identidad visual de la sede.}
  \RF{RF-44}{El envío de correo debe ejecutarse de forma asíncrona; los fallos del proveedor no deben revertir una reserva válida y deben permitir reintentos de hasta tres veces, con registro del resultado.}
  \RF{RF-45}{El sistema debe registrar el consentimiento expreso para el tratamiento de datos personales, incluyendo fecha, hora, IP y versión de la política aceptada.}
  \RF{RF-46}{El sistema debe impedir finalizar la reserva o el check-in cuando el consentimiento requerido no esté registrado.}
  \RF{RF-47}{El Administrador Global debe poder ejecutar la anonimización irreversible de los datos identificables de un huésped cuando exista una solicitud aprobada y conservar los valores contables de folios y comprobantes.}
  \RF{RF-48}{El portal de autogestión del huésped se considera una funcionalidad de fase posterior y no es obligatorio para la versión inicial.}
\end{description}

\subsection{Gerencia, reportes y auditoría}
\begin{description}[leftmargin=2.1cm,style=nextline]
  \RF{RF-49}{El sistema debe presentar un dashboard consolidado de la cadena y permitir filtrar las métricas por sede.}
  \RF{RF-50}{El dashboard debe presentar ingresos diarios, semanales y mensuales, ocupación, ADR, RevPAR y cancelaciones.}
  \RF{RF-51}{El dashboard debe actualizar las métricas sin requerir una recarga completa de la página, mediante sondeo periódico o tecnología equivalente.}
  \RF{RF-52}{El sistema debe permitir exportar transacciones y listas de huéspedes a Excel (.xlsx) y PDF según los permisos del usuario.}
  \RF{RF-53}{Los archivos exportados deben incluir formato profesional, totales y metadatos de fecha, hora y usuario emisor.}
  \RF{RF-54}{El sistema debe permitir consultar auditoría por usuario, entidad y fecha, mostrando operación y valores anteriores y nuevos.}
  \RF{RF-55}{La auditoría debe contemplar los niveles de negocio, técnico y HTTP, incluyendo operaciones, cambios y datos técnicos de las peticiones.}
\end{description}

\section{Requisitos no funcionales}

Los requisitos no funcionales establecen las condiciones de calidad, seguridad, rendimiento y operación del sistema. Los valores siguientes mantienen los objetivos medibles indicados en el documento fuente.

\begin{longtable}{p{2.5cm} p{7.3cm} p{3.6cm}}
\toprule
\textbf{ID / categoría} & \textbf{Requisito} & \textbf{Verificación}\\
\midrule
RNF-01 Rendimiento & La consulta de disponibilidad debe responder con un percentil 95 (p95) inferior a 800 ms bajo el escenario de prueba acordado. & Prueba de carga con datos semilla y medición de latencias.\\
RNF-02 Concurrencia & El sistema debe soportar 50 usuarios simultáneos sin errores de integridad de datos. & Prueba de carga y prueba concurrente de reservas.\\
RNF-03 Disponibilidad & El objetivo de disponibilidad es del 99\% durante el horario operativo; el horario y la ventana de medición deben acordarse. & Health checks y reporte de despliegue.\\
RNF-04 Respaldo y recuperación & La base de datos debe contar con respaldo diario y un procedimiento de restauración documentado y probado. & Ejecución demostrable del script y prueba de restauración.\\
RNF-05 Seguridad & La revisión no debe identificar vulnerabilidades críticas del OWASP Top 10 en el alcance evaluado. & Lista de verificación OWASP y análisis estático.\\
RNF-06 Accesibilidad & Las pantallas críticas de reserva, check-in y folio deben cumplir WCAG 2.1 nivel AA. & Revisión con Lighthouse/axe y navegación por teclado.\\
RNF-07 Compatibilidad & La interfaz debe funcionar en las dos últimas versiones de Chrome, Edge y Firefox, en escritorio y tableta. & Pruebas manuales documentadas en navegadores y dispositivos objetivo.\\
RNF-08 Mantenibilidad & La cobertura de pruebas debe ser al menos del 80\% en las capas Domain y Application; el quality gate de SonarQube debe aprobarse. & Informe de cobertura en CI y resultado de SonarQube.\\
RNF-09 Observabilidad & Los registros deben ser estructurados e incluir un identificador de correlación por petición. & Inspección de logs durante una demostración.\\
RNF-10 Aislamiento multitenant & Toda consulta o modificación operativa debe respetar el hotel asociado al contexto activo, salvo la consulta de disponibilidad entre sedes expresamente autorizada. & Pruebas de integración por módulo intentando acceder a datos de otra sede.\\
RNF-11 Integridad de reservas & La prevención del overbooking debe garantizarse en la base de datos y no únicamente mediante validaciones en código. & Pruebas de reservas simultáneas y revisión de restricciones/índices.\\
RNF-12 Seguridad de sesión & Los tokens de acceso deben durar entre 15 y 30 minutos. Los refresh tokens deben rotarse, almacenarse de forma segura y entregarse mediante cookies HttpOnly, Secure y SameSite=Strict en navegador. & Pruebas de autenticación, expiración, revocación y configuración de cookies.\\
RNF-13 Protección de contraseñas & Las contraseñas deben almacenarse usando BCrypt con costo adecuado o Argon2id; nunca deben guardarse en texto plano ni incluirse en logs. & Revisión de implementación y configuración.\\
RNF-14 Interfaz adaptable & Las pantallas deben ser responsive y adecuadas para resoluciones de escritorio de 1920x1080 y 1366x768, además de tabletas. & Pruebas visuales y funcionales en resoluciones objetivo.\\
RNF-15 Usabilidad y validación & Los formularios deben mostrar errores específicos, deshabilitar el guardado cuando sean inválidos y proporcionar mensajes accionables. Las tablas deben paginar en servidor y aplicar debounce de 300 ms en búsquedas. & Pruebas funcionales y revisión de interfaz.\\
RNF-16 Trazabilidad & Las operaciones sensibles deben conservar registro de usuario, fecha y cambios, respetando las políticas de borrado lógico y auditoría. & Pruebas de auditoría y revisión de registros.\\
\bottomrule
\end{longtable}

\section{Historias de usuario}

Las historias de usuario conservan el formato \textit{Como [rol], quiero [necesidad], para [beneficio]}. Las prioridades y fases se toman de la especificación del proyecto. Los criterios de aceptación se expresan en formato Dado/Cuando/Entonces y se resumen para facilitar su trazabilidad.

\subsection{Épica 0: Acceso, usuarios y seguridad}

\HU{HU-01}{Inicio de sesión y sesión segura}{Must}{1}{Como usuario del sistema, quiero iniciar sesión con correo y contraseña y mantener una sesión segura, para acceder únicamente a las funciones permitidas por mi rol.}
\begin{itemize}
  \item Dado que las credenciales son válidas, cuando inicio sesión, entonces recibo un token de acceso de 15 a 30 minutos y un refresh token rotativo seguro.
  \item Dado que se intenta reutilizar un refresh token usado o revocado, cuando se solicita renovar la sesión, entonces se rechaza y se revoca la cadena de tokens.
  \item Dado que ocurren cinco intentos fallidos consecutivos, cuando se intenta un sexto, entonces la cuenta se bloquea temporalmente y el evento queda auditado.
\end{itemize}

\HU{HU-02}{Empleados, roles y asignación a hoteles}{Must}{1}{Como Administrador Global, quiero crear empleados con rol, departamento, turno y hoteles asignados, para controlar quién opera en cada sede.}
\begin{itemize}
  \item Dado un empleado asignado a un hotel, cuando inicia sesión, entonces queda limitado al contexto de ese hotel.
  \item Dado un empleado asignado a varios hoteles, cuando inicia sesión, entonces puede elegir y cambiar únicamente entre sus sedes asignadas.
  \item Dado un empleado desactivado, cuando se confirma la desactivación, entonces se revocan sus refresh tokens y se conserva su historial mediante borrado lógico.
\end{itemize}

\subsection{Épica 1: Parametrización de la cadena}

\HU{HU-03}{Gestión de hoteles}{Must}{1}{Como Administrador Global, quiero crear y configurar sedes con nombre, ciudad, dirección, zona horaria, moneda e identidad visual, para administrar la cadena sin depender del equipo de desarrollo.}
\begin{itemize}
  \item Dado un hotel con datos válidos, cuando lo guardo, entonces queda disponible para asignar empleados y habitaciones.
  \item Dado un hotel con reservas futuras activas, cuando intento desactivarlo, entonces la operación se rechaza e informa cuántas reservas la impiden.
\end{itemize}

\HU{HU-04}{Tipos de habitación y habitaciones}{Must}{1}{Como Administrador Global o de Sede, quiero definir tipos de habitación, capacidades, comodidades y habitaciones físicas, para mantener parametrizado el inventario real de cada hotel.}
\begin{itemize}
  \item Dado un número de habitación ya registrado en el mismo hotel, cuando intento crearlo nuevamente, entonces el sistema lo rechaza; el mismo número puede existir en otra sede.
  \item Dada una habitación con reservas futuras activas, cuando intento desactivarla, entonces el sistema bloquea la operación y explica el motivo.
\end{itemize}

\HU{HU-05}{Temporadas, fines de semana y promociones}{Must}{1}{Como Administrador Global, quiero definir tarifas base, temporadas, recargos de fin de semana y promociones, para calcular precios de manera consistente.}
\begin{itemize}
  \item Dadas dos temporadas con fechas traslapadas para una misma configuración, cuando intento guardarlas, entonces el sistema rechaza el traslape.
  \item Dada una promoción vigente, cuando una reserva cae dentro de su vigencia, entonces el descuento se aplica y aparece en el desglose.
\end{itemize}

\HU{HU-06}{Impuestos y políticas de cancelación por sede}{Must}{1}{Como Administrador Global, quiero asignar impuestos y políticas de cancelación a cada sede, para respetar las reglas comerciales y normativas correspondientes.}
\begin{itemize}
  \item Dado un impuesto configurado, cuando se calcula una reserva, entonces se agrega por separado al folio.
  \item Dada una política de cancelación, cuando se cancela una reserva, entonces la penalidad se calcula con la política vigente al momento de reservar.
\end{itemize}

\HU{HU-07}{Servicios adicionales}{Should}{1}{Como Administrador de Sede, quiero administrar el catálogo de servicios y sus precios por hotel, para cargar correctamente los consumos al folio del huésped.}
\begin{itemize}
  \item Dado un servicio inactivo, cuando recepción registra un consumo, entonces el servicio no aparece como opción seleccionable.
\end{itemize}

\subsection{Épica 2: Disponibilidad y reservas}

\HU{HU-08}{Consulta de disponibilidad y cotización}{Must}{1}{Como Recepcionista, quiero consultar habitaciones disponibles por fechas, tipo y número de huéspedes con el precio total, para responder al cliente en pocos segundos.}
\begin{itemize}
  \item Dado un rango de fechas, cuando consulto, entonces veo las habitaciones libres de mi hotel y el precio total desglosado por noche.
  \item Dada una habitación con una reserva activa que cubre parte del rango, cuando consulto, entonces no aparece como disponible.
\end{itemize}

\HU{HU-09}{Disponibilidad entre sedes}{Must}{2}{Como Recepcionista, quiero consultar disponibilidad en otras sedes, para ofrecer alternativas sin realizar llamadas telefónicas.}
\begin{itemize}
  \item Dada una consulta, cuando se ejecuta, entonces se muestran solo sede, tipo de habitación, cantidad disponible y tarifa referencial.
  \item Dada cualquier respuesta, entonces no se exponen reservas, huéspedes, números de habitación ni consumos de otras sedes.
  \item Dada una consulta entre sedes, entonces su uso queda registrado en la auditoría HTTP.
\end{itemize}

\HU{HU-10}{Crear reserva con cálculo de tarifa}{Must}{1}{Como Recepcionista, quiero crear una reserva en un único flujo con huésped titular, fechas y habitación, para reducir pasos manuales y errores.}
\begin{itemize}
  \item Dada una reserva válida, cuando la guardo, entonces se crea en estado Pendiente con tarifa, descuentos e impuestos desglosados.
  \item Dada la misma petición con la misma clave de idempotencia, cuando se repite, entonces se devuelve la respuesta original y no se crea otra reserva.
\end{itemize}

\HU{HU-11}{Prevención de overbooking}{Must}{1}{Como Recepcionista, quiero que el sistema impida asignar una habitación a dos reservas con fechas traslapadas, para evitar sobreventas.}
\begin{itemize}
  \item Dadas dos peticiones simultáneas para la misma habitación y fechas, cuando se procesan, entonces solo una operación puede reservarla y la otra recibe un conflicto 409.
  \item Dado que una reserva termina el mismo día en que otra comienza, entonces no se considera traslape.
  \item Dada una edición concurrente de una misma reserva, cuando dos usuarios guardan cambios, entonces se detecta el conflicto y se responde 409.
\end{itemize}

\HU{HU-12}{Máquina de estados de la reserva}{Must}{1}{Como Recepcionista, quiero cambiar el estado de una reserva según el flujo definido, para garantizar la trazabilidad de la estancia.}
\begin{itemize}
  \item Dada una transición no permitida, cuando se intenta, entonces se rechaza con error 422.
  \item Dada una transición válida, entonces se registra en el historial el usuario, la fecha, el estado anterior, el nuevo estado y el motivo cuando corresponda.
\end{itemize}

\HU{HU-13}{Cancelación, No-Show y reembolso}{Must}{2}{Como Recepcionista o Administrador de Sede, quiero anular reservas o marcar No-Show aplicando la política del hotel, para calcular correctamente penalidades y liberar la habitación.}
\begin{itemize}
  \item Dada una cancelación dentro de la ventana gratuita, cuando se confirma, entonces la penalidad es cero y se genera el comprobante de saldo a favor si hubo abono.
  \item Dada una cancelación fuera de la ventana gratuita, cuando se confirma, entonces se calcula y muestra la penalidad antes de finalizar.
  \item Dada la confirmación de la cancelación, entonces la habitación queda disponible para esas fechas.
\end{itemize}

\HU{HU-14}{Calendario de ocupación}{Should}{2}{Como Recepcionista, quiero consultar un calendario de habitaciones por día con sus reservas, para conocer rápidamente el estado de ocupación.}
\begin{itemize}
  \item Dado un hotel y un rango de 14 días, cuando abro el calendario, entonces veo las habitaciones y sus reservas con estados diferenciados por color y texto.
\end{itemize}

\subsection{Épica 3: Operación en recepción}

\HU{HU-15}{Check-in y registro de huéspedes}{Must}{1}{Como Recepcionista, quiero registrar al titular y sus acompañantes y realizar el check-in, para cumplir el protocolo de ingreso.}
\begin{itemize}
  \item Dada una habitación que no está Limpia/Disponible, cuando intento el check-in, entonces se bloquea la operación y se sugiere una alternativa.
  \item Dado un check-in completado, entonces la reserva pasa a Check-In y la habitación a Ocupada.
  \item Dado que no existe consentimiento de tratamiento de datos, entonces no se permite completar el check-in.
\end{itemize}

\HU{HU-16}{Housekeeping}{Must}{2}{Como Camarera o Supervisor, quiero consultar y cambiar desde una tableta el estado de limpieza de mis habitaciones, para que recepción conozca en tiempo real cuáles están listas.}
\begin{itemize}
  \item Dado un check-out completado, entonces la habitación pasa automáticamente a Sucia.
  \item Dada una habitación cuya limpieza terminó, cuando la camarera la marca como Limpia, entonces recepción ve el estado actualizado sin recargar.
  \item Dada una tableta de 10 pulgadas, entonces los controles táctiles son suficientemente amplios para usarse cómodamente.
\end{itemize}

\HU{HU-17}{Consumos en habitación y folio}{Must}{1}{Como Recepcionista, quiero registrar cargos de minibar, restaurante, lavandería y otros servicios en el folio, para evitar cobros omitidos.}
\begin{itemize}
  \item Dado un servicio del catálogo, cuando registro un consumo, entonces queda en el folio con cantidad, valor, usuario y hora.
  \item Dado un cargo que requiere corrección, entonces no se edita directamente: se anula con motivo mediante un contra-asiento auditado.
  \item Dado un folio consultado, entonces se muestran cargos, impuestos, pagos y saldo.
\end{itemize}

\HU{HU-18}{Pagos y abonos}{Must}{1}{Como Recepcionista, quiero registrar pagos y abonos con método y referencia, para mantener actualizado el saldo del folio.}
\begin{itemize}
  \item Dado un pago con método y referencia, cuando lo registro, entonces reduce el saldo y genera un recibo.
  \item Dada la misma petición con la misma clave de idempotencia, cuando se repite, entonces el pago no se duplica.
\end{itemize}

\HU{HU-19}{Check-out y comprobante}{Must}{1}{Como Recepcionista, quiero cerrar la estancia y emitir el comprobante final, para liquidar la cuenta de forma rápida y correcta.}
\begin{itemize}
  \item Dado un saldo distinto de cero, cuando intento el check-out, entonces se bloquea la operación y se informa el valor pendiente.
  \item Dado un saldo igual a cero, cuando realizo el check-out, entonces la reserva pasa a Check-Out, la habitación a Sucia y se genera un PDF.
  \item Dado el folio cerrado y el comprobante final emitido, entonces la reserva pasa a Liquidada.
\end{itemize}

\HU{HU-20}{Cierre de caja por turno}{Should}{2}{Como Recepcionista o Administrador de Sede, quiero consultar y cerrar el total de pagos recibidos en mi turno por método de pago, para cuadrar caja sin hojas de cálculo.}
\begin{itemize}
  \item Dado un turno, cuando lo cierro, entonces se congelan los totales por método de pago y el cierre no admite nuevos pagos.
\end{itemize}

\subsection{Épica 4: Experiencia y privacidad del huésped}

\HU{HU-21}{Notificaciones por correo}{Must}{2}{Como Huésped, quiero recibir correos automáticos con el detalle de mi reserva, voucher con código QR y comprobantes, para disponer de respaldo digital inmediato.}
\begin{itemize}
  \item Dada una reserva confirmada, entonces se envía un correo de forma asíncrona sin retrasar la respuesta HTTP.
  \item Dado un fallo del proveedor de correo, entonces la reserva no se afecta, el envío se reintenta hasta tres veces y el resultado queda registrado.
  \item Dada la plantilla del correo, entonces es adaptable y usa la identidad visual del hotel correspondiente.
\end{itemize}

\HU{HU-22}{Consentimiento de datos (Habeas Data)}{Must}{2}{Como Huésped, quiero autorizar expresamente el tratamiento de mis datos personales, para que se respeten mis derechos de privacidad.}
\begin{itemize}
  \item Dada una reserva o un check-in, cuando falta el consentimiento, entonces no se puede finalizar el proceso.
  \item Dado el consentimiento, entonces se guardan fecha, hora, IP y versión de la política aceptada.
\end{itemize}

\HU{HU-23}{Derecho al olvido y anonimización}{Should}{3}{Como Administrador Global, quiero anonimizar los datos identificables de un huésped cuando exista una solicitud aprobada, para atender la solicitud sin perder la información contable requerida.}
\begin{itemize}
  \item Dada una solicitud aprobada, cuando se ejecuta, entonces nombre, documento, correo y teléfono se reemplazan por valores irreversibles.
  \item Dada la anonimización, entonces los valores numéricos de folios y comprobantes permanecen intactos.
  \item Dada la operación, entonces queda auditada y no puede deshacerse.
\end{itemize}

\HU{HU-24}{Portal del huésped}{Could}{3}{Como Huésped, quiero consultar y reservar en línea, para autogestionar mi estancia.}
\begin{itemize}
  \item Dado un huésped autenticado, cuando consulta sus reservas, entonces solo puede ver las que le pertenecen.
  \item Esta funcionalidad es de fase posterior y su alcance debe confirmarse según el tiempo disponible.
\end{itemize}

\subsection{Épica 5: Gerencia, reportes y auditoría}

\HU{HU-25}{Dashboard consolidado multihotel}{Must}{2}{Como Socio, quiero consultar indicadores casi en tiempo real de la cadena y de cada sede, para tomar decisiones con información actualizada.}
\begin{itemize}
  \item Dado el dashboard, cuando selecciono la cadena, entonces veo ingresos diarios, semanales y mensuales, ocupación, ADR, RevPAR y cancelaciones consolidadas.
  \item Dado el selector de hotel, cuando elijo una sede, entonces todas las métricas se filtran a esa sede.
  \item Dados nuevos datos, entonces el panel se actualiza sin recargar la página completa.
\end{itemize}

\HU{HU-26}{Exportación a Excel y PDF}{Should}{2}{Como Auditor o Administrador de Sede, quiero exportar transacciones y listas de huéspedes a XLSX y PDF, para realizar conciliaciones contables.}
\begin{itemize}
  \item Dada una exportación grande, cuando se genera, entonces se procesa por streaming sin cargar todo el contenido en memoria.
  \item Dado el archivo generado, entonces incluye formato profesional, totales y metadatos de fecha, hora y usuario emisor.
\end{itemize}

\HU{HU-27}{Consulta de auditoría}{Should}{3}{Como Auditor, quiero consultar cambios por usuario, entidad y fecha, para detectar modificaciones no autorizadas o inconsistencias.}
\begin{itemize}
  \item Dado un filtro de consulta, entonces veo la operación, el usuario y los valores anteriores y nuevos.
  \item Dado un usuario con rol Auditor, entonces el acceso es de solo lectura; el Administrador de Sede solo puede consultar la auditoría de su propio hotel.
\end{itemize}

\section{Reglas de negocio transversales}

\begin{enumerate}
  \item \textbf{Aislamiento por sede:} toda entidad operativa debe asociarse a un hotel y filtrarse automáticamente por el contexto de sesión.
  \item \textbf{Disponibilidad entre sedes:} es una excepción de solo lectura y únicamente puede revelar sede, tipo de habitación, cantidad disponible y tarifa referencial.
  \item \textbf{Estados de reserva:} las transiciones no definidas deben rechazarse con error 422.
  \item \textbf{Check-in:} requiere fecha válida, habitación Limpia/Disponible y consentimiento registrado.
  \item \textbf{Check-out:} exige saldo del folio igual a cero.
  \item \textbf{Folio:} los cargos registrados son inmutables; las correcciones se realizan mediante contra-asientos auditados.
  \item \textbf{Idempotencia:} repetir una petición de creación de reserva o registro de pago con la misma clave no debe producir efectos duplicados.
  \item \textbf{Concurrencia:} el control de versión de fila ayuda a detectar ediciones simultáneas, pero la prevención de overbooking requiere una restricción de integridad en la base de datos.
  \item \textbf{Pagos en línea:} no se incluye una pasarela de pagos en esta versión; los pagos se registran manualmente.
  \item \textbf{Privacidad:} solo se deben recopilar datos necesarios para la estancia y debe conservarse evidencia del consentimiento.
\end{enumerate}

\section{Trazabilidad entre requisitos e historias de usuario}

\begin{longtable}{p{4.1cm} p{4cm} p{5.2cm}}
\toprule
\textbf{Área} & \textbf{Requisitos relacionados} & \textbf{Historias relacionadas}\\
\midrule
Acceso y seguridad & RF-01 a RF-07; RNF-12, RNF-13 & HU-01, HU-02\\
Parametrización & RF-08 a RF-17 & HU-03, HU-04, HU-05, HU-06, HU-07\\
Disponibilidad y reservas & RF-18 a RF-29; RNF-01, RNF-02, RNF-10, RNF-11 & HU-08 a HU-14\\
Recepción y estancia & RF-30 a RF-42 & HU-15 a HU-20\\
Notificaciones y privacidad & RF-43 a RF-48 & HU-21 a HU-24\\
Gerencia y auditoría & RF-49 a RF-55; RNF-09, RNF-16 & HU-25, HU-26, HU-27\\
Calidad transversal & RNF-03 a RNF-08, RNF-14, RNF-15 & Aplica a todas las historias según la pantalla o proceso involucrado\\
\bottomrule
\end{longtable}

\section{Supuestos y aspectos por confirmar}

Los siguientes puntos se mantienen como decisiones o asuntos pendientes del documento fuente y no deben considerarse cerrados sin validación:
\begin{itemize}
  \item Confirmar con la Junta la definición operativa de \textit{Liquidada}: folio cerrado y comprobante final emitido.
  \item Ajustar y acordar el horario operativo para medir el objetivo de disponibilidad del 99\%.
  \item Confirmar las obligaciones legales y contables aplicables a la conservación de folios y comprobantes después de una anonimización.
  \item La multimoneda es deseable, pero no obligatoria; la moneda base prevista es COP.
  \item La matriz de roles y permisos es una propuesta que debe validar la Junta.
  \item El portal de huésped corresponde a una fase posterior.
  \item El documento fuente menciona SQL Server como base preferida y también describe alternativas técnicas para imponer la no superposición de fechas; la implementación concreta debe verificarse con el motor finalmente seleccionado.
\end{itemize}

\section{Conclusión}

Los requisitos funcionales describen las operaciones que debe ofrecer el sistema, desde el acceso y la configuración de sedes hasta las reservas, la estancia, los pagos y los reportes. Los requisitos no funcionales establecen objetivos medibles de rendimiento, disponibilidad, seguridad, accesibilidad, compatibilidad, mantenibilidad y observabilidad. Las 27 historias de usuario conectan esas capacidades con las necesidades de los actores del sistema y aportan criterios de aceptación que pueden utilizarse para planificar el desarrollo y diseñar las pruebas.

\end{document}
